% ════════════════════════════════════════════════════════════
% Chapter 2 — Technology Foundations
% Video-to-Gaussian: An Agentic Pipeline for Volumetric
% Cultural Heritage Preservation (v2.0)
%
% Dr John O'Hare, DreamLab AI Consulting Ltd
% For: University of Salford, June 2026
% ════════════════════════════════════════════════════════════

\chapter{Technology Foundations}\label{ch:foundations}

This chapter establishes the theoretical and algorithmic bedrock on which the
\texttt{gaussian-toolkit} pipeline is constructed. The six sections below
progress from classical photogrammetric scene recovery to the explicit volumetric
primitives, surface-extraction families, video-consistent segmentation, generative
completion, and scene-description standards that together enable a single 4K
gallery walkthrough to become a hierarchical, editable, \acrshort{usd}-encoded
volumetric record.

% ────────────────────────────────────────────────────────────
\section{From Photogrammetry to Radiance Fields}\label{sec:sfm-to-nerf}
% ────────────────────────────────────────────────────────────

Classical multi-view reconstruction decomposes into two sequentially coupled
problems: recovering camera geometry, and then recovering scene geometry.

\paragraph{Feature matching and Structure from Motion.}
\textcite{Lowe2004} demonstrated that scale-invariant feature keypoints can be
matched reliably across wide-baseline image pairs despite changes in illumination
and viewpoint. \acrshort{sfm} solves the camera-recovery problem by jointly
estimating intrinsic and extrinsic parameters from such matches through
incremental bundle adjustment. \textcite{Schoenberger2016sfm} systematised this
into \acrshort{colmap}, the de-facto reference implementation that
\texttt{gaussian-toolkit} invokes for pose estimation. Given $N$ frames, \acrshort{sfm}
yields a sparse point cloud and a set of calibrated cameras---sufficient to
anchor further reconstruction, but too sparse to serve as a visual medium.

\paragraph{Multi-View Stereo.}
Dense scene geometry is recovered by \acrlong{mvs}
\parencite{Schoenberger2016mvs}, which sweeps depth hypotheses per pixel and
aggregates photoconsistency evidence across views. The result is a dense point
cloud or depth map per image from which a fused mesh or volume can be derived.
\acrshort{mvs} reaches geometric fidelity adequate for survey applications, but
its outputs require manual clean-up for experiential rendering, and its Lambertian
assumptions break on specular, transparent, or fine-structured surfaces---all of
which occur in gallery collections.

\paragraph{Neural Radiance Fields.}
\textcite{Mildenhall2020} introduced Neural Radiance Fields (\acrshort{nerf}),
representing a scene as a continuous, differentiable function
$f_\theta(\mathbf{x}, \mathbf{d}) \to (\mathbf{c}, \sigma)$ mapping a 3D position
$\mathbf{x}$ and view direction $\mathbf{d}$ to emitted colour $\mathbf{c}$ and
volumetric density $\sigma$, parameterised by a multi-layer perceptron. Novel views
are rendered by integrating colour along rays via classical volume rendering, and
the network is optimised end-to-end against a set of posed photographs.
\acrshort{nerf} surpassed classical methods on view-synthesis quality, particularly
for reflective and translucent materials, and its implicit representation implicitly
handles the occlusion ambiguities that degrade \acrshort{mvs} dense-depth estimation.

\textcite{Barron2022MipNeRF360} extended the formulation to unbounded outdoor and
indoor scenes via a contracted coordinate system and an integrated positional
encoding that resolves the aliasing artefacts present in the original model.
Mip-\acrshort{nerf}~360 remains the quality reference against which splatting
methods are benchmarked on indoor scenes comparable to those encountered in
gallery preservation.

\paragraph{Why neural volumetric representations win for experiential interiors.}
Gallery interiors present a combination of continuous smooth wall surfaces, fine
decorative detail, mixed artificial and natural illumination, and objects whose
rear faces are never directly observed during a linear walkthrough. Classical
photogrammetry either fails on textureless or specular surfaces, or requires
controlled multi-direction captures incompatible with public-facing collection
environments. Neural volumetric methods, trained entirely from an uncontrolled
monocular video, produce renderings that preserve perceptual continuity of
illumination and material across novel viewpoints. This experiential continuity
is the primary preservation goal---the scene must be \emph{inhabitable}, not merely
measurable.

% ────────────────────────────────────────────────────────────
\section{3D Gaussian Splatting}\label{sec:3dgs}
% ────────────────────────────────────────────────────────────

\acrlong{3dgs} (\acrshort{3dgs}) \parencite{Kerbl2023} replaces the implicit neural
representation of \acrshort{nerf} with a collection of explicit, anisotropic
Gaussian primitives that can be rasterised in real time without ray marching.

\paragraph{The Gaussian primitive.}
Each primitive is characterised by a mean position $\boldsymbol{\mu} \in \mathbb{R}^3$,
an opacity $\alpha \in [0,1]$, a set of spherical-harmonic coefficients encoding
view-dependent colour, and a full covariance matrix $\boldsymbol{\Sigma}$. The
spatial contribution of a primitive at position $\mathbf{x}$ is:

\begin{equation}
  G(\mathbf{x}) = \exp\!\left(-\tfrac{1}{2}(\mathbf{x} - \boldsymbol{\mu})^\top
  \boldsymbol{\Sigma}^{-1}(\mathbf{x} - \boldsymbol{\mu})\right).
  \label{eq:gaussian}
\end{equation}

To guarantee that $\boldsymbol{\Sigma}$ is symmetric positive semi-definite throughout
optimisation, it is factored as

\begin{equation}
  \boldsymbol{\Sigma} = \mathbf{R}\,\mathbf{S}\,\mathbf{S}^\top\mathbf{R}^\top,
  \label{eq:covariance}
\end{equation}

where $\mathbf{R}$ is the rotation matrix derived from a learnt unit quaternion and
$\mathbf{S} = \mathrm{diag}(s_x, s_y, s_z)$ encodes the three principal-axis
scales. Gradients with respect to $\mathbf{R}$ and $\mathbf{S}$ are propagated
through this factorisation during back-propagation, ensuring the covariance
remains valid.

\paragraph{Differentiable rasterisation.}
Rendering follows the splatting paradigm of \textcite{Zwicker2001}: each primitive
is projected into screen space as a 2D Gaussian via an affine approximation of
the projective transformation, sorted by depth, and composited front-to-back.
The accumulated colour at a pixel is:

\begin{equation}
  C = \sum_{i=1}^{N} \mathbf{c}_i\,\alpha_i
  \prod_{j=1}^{i-1}\!\left(1 - \alpha_j\right),
  \label{eq:blending}
\end{equation}

where $\mathbf{c}_i$ is the view-dependent colour of the $i$-th primitive and the
product term models accumulated transmittance. This closed-form alpha-blending is
differentiable end-to-end, enabling photometric loss to be back-propagated through
rasterisation to all primitive parameters. The \acrshort{gpu} tile-based rasteriser
delivers interactive frame rates (typically $>$60\,fps at 1080p) without neural
network inference at render time---a decisive advantage over \acrshort{nerf} for
web delivery.

\paragraph{Densification and optimisation.}
The canonical Kerbl et al.\ optimiser employs adaptive density control: under-
reconstructed regions where the view-space positional gradient exceeds a threshold
are addressed by cloning or splitting Gaussians, while near-transparent primitives
are periodically pruned. This heuristic works well on static, well-sampled scenes
but can over-densify cluttered regions and under-densify large homogeneous surfaces.

\paragraph{Refinements in use.}
\textcite{Yu2024MipSplatting} identify a mismatch between 3D Gaussian extent and
2D rasterised footprint at different scales, introducing a 3D smoothing filter
and a 2D Mip filter that together eliminate the floaters and high-frequency
aliasing visible at oblique or close viewpoints. Mip-Splatting is the default
rendering path in \texttt{gaussian-toolkit}'s viewer module.

\textcite{Kheradmand2024} reframe densification as Markov Chain Monte Carlo
(\acrshort{mcmc}) sampling of a stochastic process over the Gaussian population:
relocation moves teleport low-utility primitives to high-gradient regions, while
birth--death moves control scene capacity. This eliminates the hand-tuned
gradient-threshold heuristics of the original densifier and produces more
uniform coverage of challenging surfaces.

\textcite{Kerbl2024Hierarchical} extend the representation to arbitrarily large
scenes by organising primitives into a spatial hierarchy of levels of detail, with
coarser levels used for distant geometry and finer levels loaded progressively as
the viewer approaches. The hierarchical formulation is directly relevant to gallery
sequences where a single video may span dozens of metres of corridor.

\paragraph{MRNF densification.}
The default densification strategy in \texttt{gaussian-toolkit} is \acrlong{mrnf}
(\acrshort{mrnf}), which combines Mini-Splatting constrained-count primitives
\parencite{Fang2024} with the \acrshort{mcmc} relocation scheme \parencite{Kheradmand2024}
and a normal-consistency regulariser that penalises primitives whose major axes
deviate from the local surface normal estimated from the sparse \acrshort{sfm}
point cloud. The normal filter suppresses the needle-shaped floating Gaussians that
tend to proliferate in poorly textured regions of gallery walls and prevents
spurious primitives from contaminating the depth maps later consumed by mesh
extractors. The net effect is a reconstruction that is both compact---typically
30--50\% fewer primitives than the canonical densifier at comparable \acrshort{psnr}---and
better conditioned for downstream meshing.

\begin{figure}[htbp]
  \centering
  \includegraphics[width=0.85\textwidth]{v2/gsplat_mesh.png}
  \caption{Co-registered Gaussian splat (left) and extracted polygonal mesh (right)
    of a gallery interior captured with \texttt{gaussian-toolkit}.  The splat
    preserves specular highlights and fine texture detail; the mesh provides
    watertight geometry suitable for physics simulation, UV unwrapping, and
    engine import. Both representations are stored as \acrshort{usd} variants of
    the same scene prim and can be toggled without reloading scene metadata.}
  \label{fig:found-splat-mesh}
\end{figure}

% ────────────────────────────────────────────────────────────
\section{Surfaces from Splats}\label{sec:splats-to-surfaces}
% ────────────────────────────────────────────────────────────

A trained \acrshort{3dgs} scene is a point process of oriented ellipsoids, not a
surface. Primitives may overlap, interpenetrate, or float in empty space. Standard
mesh operations---UV unwrapping, Boolean difference, physics collision, real-time
\acrlong{pbr} shading---require a manifold surface representation. The extraction
of such a surface from a Gaussian splat field is an active research problem, and
the \texttt{gaussian-toolkit} exposes four distinct extraction backends to allow
trade-offs between geometric fidelity, computational cost, and output topology to
be made per-project.

\paragraph{The fundamental challenge.}
A Gaussian primitive contributes density continuously through space. Unlike a
neural \acrshort{tsdf} or an occupancy network, there is no obvious iso-surface.
Na\"{i}ve density thresholding produces ragged, topologically inconsistent meshes.
Extracting geometry that faithfully captures both large smooth walls and sub-centimetre
decorative relief from the same primitive set requires principled coupling between
the volumetric representation and the surface prior.

\paragraph{Surface-aligned methods.}
\textcite{Guedon2024Sugar} introduce SuGaR, which adds a regularisation term that
encourages each Gaussian to align its flattest axis with the nearest mesh surface
during training. A refined mesh is extracted by applying the Poisson surface
reconstruction algorithm to the aligned normals. SuGaR produces clean, manifold
meshes but requires a modified training loop; it is not applicable post-hoc to a
pre-trained splat.

\textcite{Guedon2025MILo} subsequently developed \acrlong{milo} (\acrshort{milo}),
which maintains a deformable mesh in the loop during splatting optimisation. The
mesh and the Gaussians are jointly optimised: the mesh constrains the positional
prior of nearby primitives while the photometric loss drives both simultaneously.
\acrshort{milo} achieves the highest surface quality in the benchmark comparison
presented in Chapter~\ref{ch:meshing} but doubles optimisation time relative to
standard training.

\paragraph{Post-hoc confidence-based extraction.}
\textcite{Radl2025CoMe} present \acrlong{come} (\acrshort{come}), which extracts
a mesh from a \emph{trained} splat without retraining by computing a per-primitive
confidence score derived from the accumulated photometric residual. High-confidence
primitives contribute surface normals to a spatially adaptive iso-surface;
low-confidence primitives are suppressed. \acrshort{come} is well-suited to the
\texttt{gaussian-toolkit} workflow because it operates as a post-processing step,
decoupling mesh quality from the choice of training configuration.

\paragraph{Gaussian wrapping.}
\textcite{Waczynska2024GaussianWrapping} approach the problem from the inverse
direction: rather than deriving a surface from density, they initialise a proxy
mesh and bind Gaussians to its faces, constraining primitive positions and
orientations to remain coherent with the mesh throughout optimisation. The
resulting \acrlong{gw} (\acrshort{gw}) representation is geometrically exact by
construction but requires a reasonable initial mesh---practically sourced from a
classical \acrshort{tsdf} reconstruction or an \acrshort{mvs} output.

\paragraph{Classical TSDF and marching cubes.}
The \acrlong{tsdf} \parencite{Curless1996} integrates per-view depth maps into a
volumetric signed-distance field on a voxel grid. The iso-surface at zero crossing
is extracted via the marching cubes algorithm \parencite{Lorensen1987}, producing
a triangle mesh whose resolution is determined by voxel spacing. In the
\texttt{gaussian-toolkit}, the \acrshort{tsdf} backend consumes depth maps rendered
from the trained splat rather than from a physical depth sensor, giving it access
to the fully optimised scene representation. Classical \acrshort{tsdf}+marching
cubes remains a competitively efficient baseline, especially for room-scale scenes
where the Gaussian normal field is sufficiently dense to produce artefact-free
depth renders.

A systematic empirical comparison of all four backends across geometry, run-time,
and perceptual quality metrics is deferred to Chapter~\ref{ch:meshing}.

% ────────────────────────────────────────────────────────────
\section{Promptable Segmentation}\label{sec:sam}
% ────────────────────────────────────────────────────────────

Scene decomposition into individually addressable objects is a prerequisite for
constructing a meaningful scene graph: a monolithic scene primitive cannot carry
per-object metadata, cannot be independently replaced with a generatively completed
mesh, and cannot be selected by a curator examining the scene in a browser viewer.

\paragraph{Segment Anything Model.}
\textcite{Kirillov2023} introduced \acrlong{sam} (\acrshort{sam}), a promptable
image segmentation foundation model trained on a dataset of over one billion masks.
Given a prompt in the form of a point, bounding box, or free-form mask, \acrshort{sam}
returns a segmentation of the corresponding object at interactive speed. The
promptable interface maps naturally onto the \texttt{gaussian-toolkit}'s interactive
web \acrshort{usd} viewer, where a curator can click an exhibit to trigger object
isolation.

\paragraph{SAM 2 for video-consistent decomposition.}
Image-level segmentation applied to individual video frames produces inconsistent
masks: the same object receives different segmentation boundaries in adjacent frames
due to motion blur, partial occlusion, and lighting change. \textcite{Ravi2024SAM2}
extend \acrshort{sam} to video with a streaming memory architecture that propagates
object-level mask tokens across frames. Given a prompt in a single frame,
\acrshort{sam}~2 tracks the corresponding object mask through the entire video
sequence, producing temporally consistent per-frame masks at up to 44 frames per
second on a single \acrshort{gpu}.

\paragraph{Lifting 2D masks to 3D consistency.}
The \texttt{gaussian-toolkit}'s segmentation module (\texttt{segment\_scene})
projects the per-frame \acrshort{sam}~2 masks back onto the trained Gaussian
primitive set via the calibrated cameras. Each primitive is assigned to the object
whose mask covers its projected centroid in the majority of frames in which it
is visible. This voting-based assignment is robust to single-frame mask failures
and produces 3D-consistent object labels without requiring any 3D supervision.
The labelled primitive subsets form the basis of the scene-graph decomposition
described in Chapter~\ref{ch:pipeline}, with each object's Gaussians isolated,
independently meshed, and associated with a distinct \acrshort{usd} prim.

% ────────────────────────────────────────────────────────────
\section{Generative 3D Reconstruction}\label{sec:generative}
% ────────────────────────────────────────────────────────────

A hand-held camera walkthrough systematically fails to observe the posterior
surface of any forward-facing exhibit. A showcase item presented against a wall
will have its back face entirely occluded throughout the capture; a sculpture on
a plinth may have its base inaccessible. Photogrammetric methods cannot reconstruct
surfaces they have never observed. The resulting holes in the mesh are either
filled by aggressive interpolation, which introduces geometrically incorrect
surface, or left open, which disqualifies the mesh from physical simulation and
game-engine import. For cultural heritage purposes, a complete object mesh---even
one whose occluded regions are generatively inferred---is of substantially greater
archival and display value than a fragment.

\paragraph{Multi-view diffusion.}
Multi-view diffusion models learn the joint distribution of object appearance
across a canonical set of viewpoints, conditioned on one or more observed views.
They encode strong shape priors from training on large 3D asset datasets and can
hallucinate plausible geometry for unobserved surfaces while remaining consistent
with the photographic evidence on the visible faces.

\textcite{Zhao2025Hunyuan} present Hunyuan3D~2.0, a two-stage system comprising
a multi-view diffusion model that synthesises consistent views at 24 canonical
camera positions and a reconstruction model that converts those views into a
textured mesh via differentiable rasterisation. Hunyuan3D~2.0 achieves state-of-the-art
quality on object-level reconstruction from a single image and scales to
high-resolution textures (up to $2048^2$ per chart).

\paragraph{ComfyUI orchestration.}
\texttt{gaussian-toolkit} executes Hunyuan3D~2.0 through \textcite{ComfyUI2024},
a modular, node-graph execution engine for diffusion pipelines. The \texttt{generate\_object\_mesh}
module constructs a \texttt{ComfyUI} workflow graph programmatically, submitting
each isolated object's best-observed view as a conditioning image and retrieving
the resulting textured mesh via the \texttt{ComfyUI} HTTP API. Using
\texttt{ComfyUI} as the execution layer decouples the \texttt{gaussian-toolkit}
from any specific model version: as Hunyuan3D or successor models are updated,
only the workflow JSON changes.

\paragraph{Background inpainting.}
Isolating an object for generative reconstruction exposes the background behind
it. Where the object partially occludes a wall or floor, the Gaussian splat
background shows a hole at those pixels. \texttt{gaussian-toolkit} employs
FLUX.1 \parencite{FLUX2024}---a rectified flow transformer trained for
high-fidelity text-conditioned image synthesis---as an inpainting model to fill
disoccluded background regions before they are incorporated into the scene's
background mesh. FLUX.1 is conditioned on the surrounding context pixels and a
brief natural-language description of the venue, generating texture-consistent
infill that avoids the conspicuous seams produced by classical patch-based methods.

\paragraph{Why generative completion is architecturally necessary.}
The decision to adopt generative completion rather than accept incomplete meshes
follows directly from the project's archival brief. A \acrshort{usd} scene intended
for long-term deposit must contain complete, watertight geometry for each registered
object. In the absence of completion, the pipeline would require a secondary capture
pass for every non-trivial exhibit---negating the single-capture design goal.
Generative completion introduces epistemic uncertainty on unseen faces, which is
acknowledged in the per-object metadata schema via a \texttt{completionConfidence}
field; downstream consumers can filter on this field if measured geometry is
preferred.

% ────────────────────────────────────────────────────────────
\section{Scene Description}\label{sec:usd}
% ────────────────────────────────────────────────────────────

The final output of \texttt{gaussian-toolkit} is not a single mesh or splat file
but a structured, hierarchical \emph{scene}---an environment mesh, a collection
of individually completed object meshes with metadata, and the original Gaussian
splat, all co-registered and jointly deliverable.

\paragraph{Universal Scene Description.}
\acrlong{usd} (\acrshort{usd}) \parencite{Pixar2019} is an open interchange
framework developed by Pixar Animation Studios and subsequently adopted across the
visual effects, games, and digital twin industries. Its core data model is a
hierarchy of typed \emph{prims} (primitives) that can carry geometry, material
bindings, camera definitions, light sources, and arbitrary user-defined metadata
attributes. Crucially, \acrshort{usd} supports \emph{composition arcs}---mechanisms
for layering, referencing, and instancing that allow large scenes to be assembled
from independently maintained asset files. For the \texttt{gaussian-toolkit}, the
relevant composition features are:

\begin{itemize}
  \item \textbf{Variant sets}: a named collection of alternative representations
    for a prim. Each scene object carries a \texttt{representation} variant set
    with two variants---\texttt{splat} (a reference to the Gaussian primitive
    subset) and \texttt{mesh} (a reference to the extracted or generatively
    completed polygonal mesh)---enabling renderer-appropriate view switching
    without duplicating scene topology.

  \item \textbf{Layer stacking}: the environment mesh, per-object assets, and
    splat data are authored in separate \texttt{.usda} or \texttt{.usdz} files
    and composed at the root stage level, so any layer can be updated in isolation
    without invalidating the others.

  \item \textbf{Custom schemas}: the \texttt{gaussian-toolkit} registers a
    lightweight \texttt{CulturalHeritageAPI} applied schema carrying fields for
    accession number, capture date, completionConfidence, and licensing, making
    the metadata intrinsically part of the geometry asset rather than a sidecar file.
\end{itemize}

\paragraph{Browser delivery.}
The \acrshort{usd} scene is the authoritative master format; browser delivery uses
derived representations. The Gaussian splat subset is exported as a \texttt{.ksplat}
file (the compressed container supported by \texttt{ksplat}, \acrshort{ksplat})
and streamed to a \acrlong{webgl} (\acrshort{webgl}) viewer based on the
open-source renderer of \textcite{Antimatter15Splat}. The emerging \acrlong{webgpu}
(\acrshort{webgpu}) standard \parencite{Khronos2024WebGPU}, which exposes
compute shaders in the browser, enables a more efficient tiled rasterisation path
that the \texttt{gaussian-toolkit} viewer activates automatically when the browser
reports \acrshort{webgpu} availability, falling back to \acrshort{webgl} otherwise.
Polygon meshes are served as compressed \texttt{glTF~2.0} with Draco-encoded
geometry, enabling the same delivery infrastructure used for standard web 3D
content.

\paragraph{Summary.}
The six technology strata described in this chapter---\acrshort{sfm}/\acrshort{mvs}
pose recovery, \acrshort{3dgs} volumetric representation, post-hoc mesh extraction,
\acrshort{sam}~2 video segmentation, multi-view generative completion, and
\acrshort{usd} scene composition---form an engineering stack in which each layer
consumes the output of its predecessor. No layer has a fundamental dependency on
the specific implementation choice of any other: the meshing backends are
interchangeable, the generative model is invoked through an abstracted
\texttt{ComfyUI} graph, and the \acrshort{usd} representation is independent of
the delivery target. This composability is the architectural principle that allows
\texttt{gaussian-toolkit} to accommodate methodological advances in any single
layer without redesigning the pipeline as a whole.
